% Originally: content/week-08.tex, \subsection{Change of variables} --- promoted to a section

\section{Change of variables}
\label{sec:change_of_variables}

\begin{theorem}[Change of variables]
\label{thm:change_of_variables}
Suppose $U,V \subseteq \mathbb{R}^n$ are open sets and $\Phi : U \to V$ is a
$C^1$-diffeomorphism. If $A \subseteq U$ is Jordan measurable with $\overline{A} \subseteq U$,
and $f : A \to \mathbb{R}$ is continuous, then
\[\int_A f(x)\,dx
  = \int_{\Phi(A)} \frac{f\circ\Phi^{-1}(y)}{\big|\det \Jac\Phi\big(\Phi^{-1}(y)\big)\big|}\,dy.\]
\end{theorem}

\begin{ainote}
Corsin uses \newterm{Jordan measurable} \germanfor{Jordan-messbar} here without defining it;
the definition is in the lecture (ch.~13). For this chapter it is enough to know that boxes,
and the bounded regions cut out by finitely many smooth curves that appear in the examples, are
all Jordan measurable.

On the shape of the formula: it is often quoted the other way round, as
$\int_{\Phi(A)} f = \int_A (f\circ\Phi)\,\lvert\det\Jac\Phi\rvert$. The version above is the same
statement with $\Phi$ pushed to the other side, which is why the determinant appears
\emph{divided} rather than multiplied. Which form is convenient depends on whether it is $\Phi$
or $\Phi^{-1}$ that you can write down. In \cref{ex:change_of_variables_domain} below it is
$\Phi$, so the dividing form is the one that saves work.

% Added 2026-08-13 (Claude Opus 5, Extra): the two halves of the proof are worked out at the very
% end of this section, six hundred lines below, and nothing here said so. A reader has no way to
% find them.
The theorem is not proved here, but the two steps that carry its proof are, and both sit at the
end of this section. \Cref{ex:linear_substitution_elementary_matrices} settles the linear case,
where the determinant enters and everything else is bookkeeping.
\Cref{ex:cube_covering_lemma_substitution} settles the geometric step that a Lipschitz estimate
cannot reach, namely that a diffeomorphism close to the identity cannot let a cube collapse.
\end{ainote}

% Source: Jérôme Paschoud/Notizen/Woche 8.2 Riemann-Integral.pdf, p. 2
% Generator: Gemini 3.1 Pro (High)
\begin{remark}[Intuition for the Jacobian determinant]
From linear algebra, we know that for a linear map given by a matrix $L$, the absolute value of its determinant $\lvert\det L\rvert$ describes how much the map scales volumes. For example, the map $L = \left(\begin{smallmatrix} 2 & 3 \\ 1 & -1 \end{smallmatrix}\right)$ expands areas by a factor of $\lvert\det L\rvert = \lvert-2 - 3\rvert = 5$.

\begin{center}
\begin{tikzpicture}[>=Latex, scale=0.8]
  % Left side: Unit square
  \draw[->, gray] (-0.5, 0) -- (2, 0);
  \draw[->, gray] (0, -0.5) -- (0, 2);
  \draw[thick, blue, fill=blue!10] (0,0) -- (1,0) -- (1,1) -- (0,1) -- cycle;
  \draw[->, thick, blue] (0,0) -- (1,0) node[below right, font=\scriptsize, text=black] {$\left(\begin{smallmatrix} 1 \\ 0 \end{smallmatrix}\right)$};
  \draw[->, thick, blue] (0,0) -- (0,1) node[above left, font=\scriptsize, text=black] {$\left(\begin{smallmatrix} 0 \\ 1 \end{smallmatrix}\right)$};
  
  % Arrow in middle
  \draw[->, thick] (2.5, 0.5) -- (4, 0.5) node[midway, above] {$L$};
  
  % Right side: Parallelogram
  \begin{scope}[shift={(6, 1)}]
    \draw[->, gray] (-0.5, 0) -- (5.5, 0);
    \draw[->, gray] (0, -2.5) -- (0, 2.5);
    
    \coordinate (O) at (0,0);
    \coordinate (V1) at (2,1);
    \coordinate (V2) at (3,-1);
    \coordinate (V3) at (5,0);
    
    \draw[thick, blue, fill=blue!10] (O) -- (V1) -- (V3) -- (V2) -- cycle;
    \draw[->, thick, blue] (O) -- (V1) node[above left, font=\scriptsize, text=black] {$\left(\begin{smallmatrix} 2 \\ 1 \end{smallmatrix}\right)$};
    \draw[->, thick, blue] (O) -- (V2) node[below left, font=\scriptsize, text=black] {$\left(\begin{smallmatrix} 3 \\ -1 \end{smallmatrix}\right)$};
    
    % Correction: Claude Opus 5 (High) --- this caption sat at (2.5,-1.5), which is where the
    % `below left' label of V2 = (3,-1) also lands, so the column vector printed on top of
    % "Area is |det L| = 5". Dropped clear of it; nothing else lies below the parallelogram.
    \node[blue, align=center, font=\small] at (2.5, -2.1) {Area is $\lvert\det L\rvert = 5$\\ times larger};
  \end{scope}
\end{tikzpicture}
\end{center}

The change of variables formula generalizes this to non-linear maps: the Jacobian determinant acts as the local volume scaling factor.
\end{remark}

% Source: Jérôme Paschoud/Notizen/Woche 8.2 Riemann-Integral.pdf, pp. 2--3
% Generator: Gemini 3.1 Pro (High)
\begin{example}[Gaussian integral via polar coordinates]
\label{ex:gaussian_integral_polar}
A classic application of the change of variables formula is computing the Gaussian integral $I := \int_{-\infty}^\infty e^{-x^2}\,dx$. We consider its square:
\[
I^2 = \left( \int_{-\infty}^\infty e^{-x^2}\,dx \right) \left( \int_{-\infty}^\infty e^{-y^2}\,dy \right) = \int_{\mathbb{R}^2} e^{-(x^2+y^2)} \mathop{d\text{vol}}(x,y).
\]
We use polar coordinates, but strictly following \cref{thm:change_of_variables}, we define the forward map $\Phi: \mathbb{R}^2 \setminus \{(0,0)\} \to (0, \infty) \times [-\pi, \pi)$ by
\[
\Phi(x,y) := \left( \sqrt{x^2+y^2}, \operatorname{atan2}(y,x) \right) =: (r, \varphi).
\]

\begin{center}
% Correction: Claude Opus 5 (Max) --- the picture exhibited no size difference at all, under a
% caption reading "not equal in size". The two source rectangles were $[0.5,1.5]\times[0,0.5]$
% and $[0.5,1.5]\times[0.5,1]$: congruent, area $0.5$ each. Their images were the sectors
% $r\in[0.5,1.5]$ over $0$--$30^\circ$ and over $30$--$60^\circ$: also congruent, area $0.5236$
% each. Stacking the two rectangles in $\varphi$ was the cause, because $\det\Jac\Phi^{-1} = r$
% depends on $r$ alone: shifting a rectangle in $\varphi$ cannot change its image's area, and
% shifting it in $r$ is the only thing that can.
%
% The two rectangles are now side by side in $r$, both $0.5 \times 0.6$, at the same
% $\varphi$-band. Their images are then sectors of the same angular width but different radii:
%   red   $\tfrac12(45^\circ-10^\circ)(0.9^2-0.4^2) = \tfrac12(0.6109)(0.65) = 0.199$,
%   blue  $\tfrac12(45^\circ-10^\circ)(2.0^2-1.5^2) = \tfrac12(0.6109)(1.75) = 0.535$,
% a factor of $2.69$, which is the ratio of the mean radii $1.75/0.65$ as $\det\Jac = r$ predicts.
\begin{tikzpicture}[>=Latex, scale=1]
  % Left side: (r, phi) plane. Two CONGRUENT rectangles, 0.5 wide and 0.6 tall, differing only
  % in where they sit along r -- so r_2 - r_1 = r_4 - r_3 = 0.5 by construction.
  \draw[->, gray] (-0.2, 0) -- (2.6, 0) node[right, font=\scriptsize, text=black] {$r$};
  \draw[->, gray] (0, -0.2) -- (0, 1.8) node[above, font=\scriptsize, text=black] {$\varphi$};

  \draw[thick, fill=red!20]  (0.4, 0.5) rectangle (0.9, 1.1);
  \draw[thick, fill=blue!20] (1.5, 0.5) rectangle (2.0, 1.1);

  \node[left, font=\scriptsize] at (0, 1.1) {$\varphi_2$};
  \node[left, font=\scriptsize] at (0, 0.5) {$\varphi_1$};
  \node[below, font=\scriptsize] at (0.4, 0) {$r_1$};
  \node[below, font=\scriptsize] at (0.9, 0) {$r_2$};
  \node[below, font=\scriptsize] at (1.5, 0) {$r_3$};
  \node[below, font=\scriptsize] at (2.0, 0) {$r_4$};
  % Anchored south WEST, not south: centred at x = 1.2 this line is about 2.6 units wide, so it
  % reached back past x = 0 and the $\varphi$ axis printed through its opening word. Rendering
  % p.278 showed it; anchoring the left edge is the only placement that cannot drift back.
  \node[font=\scriptsize, anchor=south west] at (0.15, 1.25) {same $\Delta r$, same $\Delta\varphi$};

  % Arrow
  \draw[->, thick] (3, 0.5) .. controls (3.3, 0.8) and (3.7, 0.8) .. (4, 0.5) node[midway, above] {$\Phi^{-1}$};

  % Right side: (x,y) plane. Same angular band for both, so only the radius differs.
  \begin{scope}[shift={(6.5, 0.5)}]
    \draw[->, gray] (-0.6, 0) -- (2.6, 0) node[right, font=\scriptsize, text=black] {$x$};
    \draw[->, gray] (0, -0.6) -- (0, 2.6) node[above, font=\scriptsize, text=black] {$y$};

    \foreach \r in {0.4, 0.9, 1.5, 2.0} { \draw[gray!45, dashed] (0,0) circle (\r); }

    % Red sector: the image of the rectangle near the origin.
    \draw[thick, fill=red!20]
      (10:0.4) -- (10:0.9) arc (10:45:0.9) -- (45:0.4) arc (45:10:0.4) -- cycle;
    % Blue sector: the image of the congruent rectangle further out.
    \draw[thick, fill=blue!20]
      (10:1.5) -- (10:2.0) arc (10:45:2.0) -- (45:1.5) arc (45:10:1.5) -- cycle;
  \end{scope}

  \node[align=center, font=\small] at (4.4, -1.7)
    {two congruent rectangles, two images of \emph{different} area\\
     $\Rightarrow$ we must compensate with the (Jacobian) determinant};
\end{tikzpicture}
\end{center}

The Jacobian matrix is
\[
\Jac\Phi(x,y) = \begin{pmatrix} \frac{x}{\sqrt{x^2+y^2}} & \frac{y}{\sqrt{x^2+y^2}} \\[1ex] \frac{-y}{x^2+y^2} & \frac{x}{x^2+y^2} \end{pmatrix}.
\]
Its determinant is
\[
\det \Jac\Phi(x,y) = \frac{x^2}{(x^2+y^2)^{3/2}} + \frac{y^2}{(x^2+y^2)^{3/2}} = \frac{1}{\sqrt{x^2+y^2}} = \frac{1}{r}.
\]
Applying the theorem (ignoring the null set $\{(0,0)\}$ where $\Phi$ is undefined), we get
\[
I^2 = \int_{\Phi(\mathbb{R}^2 \setminus \{(0,0)\})} \frac{e^{-r^2}}{\lvert1/r\rvert} \,dr\,d\varphi = \int_{-\pi}^\pi d\varphi \int_0^\infty r e^{-r^2}\,dr = 2\pi \cdot \left[ -\frac{1}{2} e^{-r^2} \right]_0^\infty = 2\pi \cdot \frac{1}{2} = \pi.
\]
Since $I > 0$, we conclude that $I = \sqrt{\pi}$.
\end{example}

% Supplement: Sascha Brack/Class Notes/Week_08_Notes_Friday.pdf, pp. 2--3
\begin{example}[Change of variables on a diamond]
\label{ex:change_of_variables_diamond}
Compute $I := \int_E e^{x+y}(x-y) \mathop{d\text{vol}}(x,y)$, where $E$ is the region bounded by the lines $y=x$, $y=-x+2$, $y=x-2$, and $y=-x$.

Writing out the inequalities, $E = \{ y \leq x \} \cap \{ y \leq -x+2 \} \cap \{ y \geq -x \} \cap \{ y \geq x-2 \}$.
Rearranging these gives $0 \leq x-y \leq 2$ and $0 \leq x+y \leq 2$.
This suggests the substitution $u := x+y$ and $v := x-y$. The new domain is the square $V = [0,2] \times [0,2]$.
The inverse map is $\Phi^{-1}(u,v) = (x,y) = \big(\frac{u+v}{2}, \frac{u-v}{2}\big)$.
Its Jacobian matrix is $\Jac\Phi^{-1}(u,v) = \begin{pmatrix} 1/2 & 1/2 \\ 1/2 & -1/2 \end{pmatrix}$, so $\big|\det\Jac\Phi^{-1}\big| = \big|-\frac{1}{4} - \frac{1}{4}\big| = \frac{1}{2}$.
By the change of variables formula:
\[I = \int_0^2 \int_0^2 e^u v \left|\det\Jac\Phi^{-1}\right| \,du\,dv = \frac{1}{2} \int_0^2 \int_0^2 e^u v \,du\,dv = \frac{1}{2} \left[ e^u \right]_0^2 \left[ \frac{v^2}{2} \right]_0^2 = \frac{1}{2} (e^2-1) \cdot 2 = e^2-1.\]
\end{example}

\begin{example}[Improper integral via polar coordinates]
\label{ex:improper_integral_polar}
Compute $\int_{\mathbb{R}^2} \frac{1}{(1+x^2+y^2)^2} \mathop{d\text{vol}}(x,y)$.
This integral is improper because the domain is unbounded. Since the integrand is non-negative and continuous, its Riemann improper integral is well-defined, and we can compute it using polar coordinates $x = r\cos\theta, y = r\sin\theta$.
The Jacobian determinant is $r$. The domain $\mathbb{R}^2$ corresponds to $r \in [0, \infty)$ and $\theta \in [0, 2\pi)$.
\[
\int_{\mathbb{R}^2} \frac{1}{(1+x^2+y^2)^2} \,dx\,dy
= \int_0^{2\pi} \int_0^\infty \frac{r}{(1+r^2)^2} \,dr\,d\theta.
\]
Substituting $u = r^2$, $du = 2r\,dr$, the inner integral is $\frac{1}{2} \int_0^\infty \frac{1}{(1+u)^2} \,du = \frac{1}{2} \left[ \frac{-1}{1+u} \right]_0^\infty = \frac{1}{2}$.
The outer integral then gives $\frac{1}{2} \int_0^{2\pi} d\theta = \pi$.
\end{example}

% Supplement: Sascha Brack/Class Notes/Week_08_Notes_Monday.pdf, pp. 8--9
\begin{example}[Volume of a sphere]
\label{ex:volume_sphere}
Let $A := \{(x,y,z) \in \mathbb{R}^3 : x^2+y^2+z^2 \leq 1\}$ be the unit ball. Given that $A$ is Jordan measurable, its volume is $\mu(A) = \int_A 1 \mathop{d\text{vol}}(x,y,z)$.
We can compute this using spherical coordinates:
\[\Phi(r,\theta,\varphi) = \begin{pmatrix} r\sin\theta\cos\varphi \\ r\sin\theta\sin\varphi \\ r\cos\theta \end{pmatrix}.\]
Up to a Jordan null set (which does not change the volume), the domain $A$ corresponds to the box $V = (0,1) \times (0,\pi) \times (-\pi,\pi)$.
The Jacobian determinant is $\det\Jac\Phi(r,\theta,\varphi) = r^2\sin\theta$, which is positive on this domain.
Applying the change of variables formula and Fubini's theorem (which allows slicing up the integral):
\begin{align*}
\mu(A) &= \int_{\Phi^{-1}(A)} \big|\det\Jac\Phi(r,\theta,\varphi)\big| \,dr\,d\theta\,d\varphi \\
&= \int_{-\pi}^\pi \int_0^\pi \int_0^1 r^2\sin\theta \,dr\,d\theta\,d\varphi \\
&= \left(\int_{-\pi}^\pi d\varphi\right) \left(\int_0^\pi \sin\theta \,d\theta\right) \left(\int_0^1 r^2 \,dr\right) \\
&= (2\pi) \cdot (2) \cdot \left(\frac{1}{3}\right) = \frac{4\pi}{3}.
\end{align*}
\end{example}

% Quelle: Linus Lüchinger/Slides/Slides-05-11.pdf, slide 3
% Generator: Claude Opus 5 (High)
\begin{remark}[Polar coordinates as slicing into shells]
\label{rem:polar_as_shells}
The factor $r^2$ in the computation above, and the $r$ in the two-dimensional case, are usually
produced by grinding out a Jacobian determinant. There is a reading that gets them without any
computation, and that generalises to $\mathbb{R}^n$ at no extra cost.

Slice $\mathbb{R}^n$ into spheres centred at the origin. The sphere of radius $r$ is an
$(n-1)$-dimensional submanifold, and it is the unit sphere scaled by $r$; since scaling by $r$
multiplies $(n-1)$-dimensional volume by $r^{n-1}$, its volume is
\[\vol_{n-1}\big(r\,\mathbb{S}^{n-1}\big) = C_n\,r^{n-1},
  \qquad C_n := \vol_{n-1}\big(\mathbb{S}^{n-1}\big).\]
A shell between radii $r$ and $r+dr$ is that sphere thickened by $dr$, so it contributes
$C_n r^{n-1}\,dr$ of $n$-dimensional volume. Adding the shells,
\[\int_{B_R(0)} f(\lvert x\rvert)\,dx = C_n\int_0^R f(r)\,r^{n-1}\,dr
  \qquad\text{for radial } f.\]

The powers now need no memorising: $n=2$ gives $r^{1}$ and $n=3$ gives $r^{2}$, matching the two
Jacobians above, and $C_2 = 2\pi$, $C_3 = 4\pi$ are the circumference of the unit circle and the
area of the unit sphere. Taking $f \equiv 1$, $n=3$, $R=1$ reproduces
$4\pi\int_0^1 r^2\,dr = 4\pi/3$ in one line.

This is a genuine shortcut only for \emph{radial} integrands, where the angular integral is a
constant that can be pulled out; when $f$ depends on the angles as well, the full change of
variables is unavoidable. But it explains where the radial power comes from, which the
determinant computation does not.
\end{remark}

% Supplement: Analysis_II_Script_v1.pdf, pp. 113--114
\begin{remark}[Reference: the standard coordinate Jacobians]
\label{rem:coordinate_jacobians_table}
The polar, cylindrical and spherical substitutions recur throughout this chapter and the next
several; collected here in one place, with $r \geq 0$ throughout and each parametrizing its
target domain up to a Jordan null set:
\begin{center}
\begin{tabular}{@{}lll@{}}
\toprule
\textbf{Coordinates} & \textbf{Transformation} & $\lvert\det\Jac\Phi\rvert$ \\
\midrule
Polar ($\mathbb{R}^2$) &
  $\Phi(r,\varphi) = (r\cos\varphi,\ r\sin\varphi)$ & $r$ \\
Cylindrical ($\mathbb{R}^3$) &
  $\Phi(r,\varphi,z) = (r\cos\varphi,\ r\sin\varphi,\ z)$ & $r$ \\
Spherical ($\mathbb{R}^3$) &
  $\Phi(r,\theta,\varphi) = (r\sin\theta\cos\varphi,\ r\sin\theta\sin\varphi,\ r\cos\theta)$
  & $r^2\sin\theta$ \\
\bottomrule
\end{tabular}
\end{center}
The angular ranges are $\varphi \in (-\pi,\pi)$ throughout, and $\theta \in (0,\pi)$
(\newterm{colatitude}, measured from the positive $z$-axis) for the spherical case. Cylindrical
coordinates are just polar coordinates in the first two variables, carried along unchanged in
$z$, which is exactly why the Jacobian is the same $r$ as in the planar case: the extra $dz$
contributes a factor of $1$. Spherical coordinates are used already in
\cref{ex:volume_sphere} above; cylindrical coordinates are the natural choice whenever a region
has rotational symmetry about an axis without being a ball, e.g.\ a solid cylinder or cone.
\end{remark}

% Source: old_exams/examFS24.pdf (21 August 2024), Wahr/Falsch-Fragen, Aufgabe 8, p. 5
% Extractor: Claude Opus 5 (High)
% Worked out here rather than set as an exercise: it is a checklist against the table immediately
% above, and its three-dimensional counterpart ex:examHS24_TF_spherical is already an exercise
% further down this file. Note the exam's argument order (theta, r), which is the reverse of the
% table's; kept as the exam has it, since the sign of the determinant depends on it.
\begin{example}[Polar coordinates, checked item by item]
\label{ex:polar_coordinates_checklist}
% Generator: Claude Opus 5 (High) --- the four statements are the examiner's; the working is
% written for these notes.
Consider $\Psi : \mathbb{R}\times[0,\infty) \to \mathbb{R}^2$,
$\Psi(\theta, r) := (r\cos\theta,\ r\sin\theta)$, and note the argument order: the angle first.
Four statements, of which two are true.
\begin{enumerate}[label=\textbf{(\alph*)}]
  \item \textbf{$\Psi$ restricted to $(0,2\pi]\times(0,\infty)$ is injective: true.} Suppose
    $\Psi(\theta_1,r_1) = \Psi(\theta_2,r_2)$. Taking Euclidean norms gives $r_1 = r_2 =: r > 0$,
    and dividing by $r$ leaves $(\cos\theta_1,\sin\theta_1) = (\cos\theta_2,\sin\theta_2)$, so
    $\theta_1 - \theta_2 \in 2\pi\mathbb{Z}$. A half-open interval of length $2\pi$ contains at
    most one representative of each such class, so $\theta_1 = \theta_2$.
  \item \textbf{$\Psi$ restricted to $(0,2\pi]\times(0,\infty)$ is surjective: false.} Its image is
    $\mathbb{R}^2\setminus\{(0,0)\}$: every non-zero point is reached, by \textbf{(a)} exactly
    once, and the origin is reached only from $r = 0$, which the restriction excludes.
  \item \textbf{$\det\Jac\Psi(\theta,r) = r^2\sin\theta\cos\theta$: false.} The columns of the
    Jacobi matrix are the partial derivatives in $\theta$ and in $r$, in that order:
    \[\Jac\Psi(\theta,r) = \begin{pmatrix} -r\sin\theta & \cos\theta \\
      r\cos\theta & \sin\theta\end{pmatrix}, \qquad
      \det\Jac\Psi(\theta,r) = -r\sin^2\theta - r\cos^2\theta = -r .\]
    The magnitude $\lvert\det\Jac\Psi\rvert = r$ is the entry in
    \cref{rem:coordinate_jacobians_table}; the minus sign appears only because the exam lists the
    angle before the radius, which swaps two columns.
  \item \textbf{$\vol_2(B_R) = 2\pi\int_0^R r\,dr$ for $B_R = \{x^2+y^2 < R^2\}$: true.} By
    \cref{thm:change_of_variables} with the weight $\lvert\det\Jac\Psi\rvert = r$, and since the
    integrand $1$ does not depend on $\theta$, the angular integral contributes a factor $2\pi$.
    The right-hand side is $2\pi\cdot R^2/2 = \pi R^2$, as it should be. This is also the $n=2$
    case of \cref{rem:polar_as_shells}.
\end{enumerate}
The proposed formula in \textbf{(c)} is worth a second look, because it is not a random wrong
answer: $r^2\sin\theta$ is the \emph{spherical} Jacobian, and $\cos\theta$ is what a misremembered
product rule adds to it. There is a cheap guard against importing one row of the table into
another. Every system in it has the form $\Phi(r,\omega) = r\,u(\omega)$ for a map $u$ into the
unit sphere, so $\Phi$ is homogeneous of degree $1$ in $r$. The radial column of $\Jac\Phi$ is then
$u(\omega)$, free of $r$, while each of the $n-1$ angular columns is $r\,\partial u/\partial\omega_i$
and so carries exactly one factor $r$. Multilinearity of the determinant pulls all $n-1$ of them
out, leaving $\lvert\det\Jac\Phi\rvert = r^{n-1}$ times a function of the angles alone. In
$\mathbb{R}^2$ that power is $r^1$, and the proposal carries $r^2$.
\exinfo{This example is taken from the exam of 21 August 2024 (Prof.\ Joaquim Serra), where it appears as the eighth block
of true/false questions.}
\end{example}

% Source: Jérôme Paschoud/Notizen/Woche 9.1 Variablenwechsel.pdf, p. 2
% Generator: Gemini 3.1 Pro (High)
\begin{example}[Volume of an ellipse]
\label{ex:volume_ellipse}
Compute the area (2-dimensional volume) of the ellipse $E := \{(x,y) \in \mathbb{R}^2 \mid (x/a)^2 + (y/b)^2 < 1\}$ for $a,b > 0$.
We use the transformation $\Phi(r,\varphi) = (ar\cos\varphi, br\sin\varphi)$. This maps the rectangle $V = (0,1) \times (0,2\pi)$ bijectively onto $E$ minus a null set (the positive $x$-axis).
The Jacobian matrix is
\[\Jac\Phi(r,\varphi) = \begin{pmatrix} a\cos\varphi & -ar\sin\varphi \\ b\sin\varphi & br\cos\varphi \end{pmatrix},\]
and its determinant is $\det\Jac\Phi(r,\varphi) = abr(\cos^2\varphi + \sin^2\varphi) = abr$.
Thus, by the change of variables formula,
\[\int_E 1 \,dx\,dy = \int_0^{2\pi} \int_0^1 abr \,dr\,d\varphi = \left(\int_0^{2\pi} ab \,d\varphi\right) \left(\int_0^1 r \,dr\right) = 2\pi ab \cdot \frac{1}{2} = \pi ab.\]
\end{example}

% Source: Jérôme Paschoud/Notizen/Woche 9.1 Variablenwechsel.pdf, p. 2
% Generator: Gemini 3.1 Pro (High)
\begin{example}[Area of a cardioid]
\label{ex:area_cardioid}
Compute the area of the region $\Omega$ given in polar coordinates by $\Omega := \{(r,\varphi) \mid \varphi \in [0, 2\pi),\ 0 \leq r \leq 1 + \cos\varphi\}$.
Using polar coordinates $\Phi(r,\varphi) = (r\cos\varphi, r\sin\varphi)$, the area is
\[\int_\Omega 1 \,dx\,dy = \int_0^{2\pi} \int_0^{1+\cos\varphi} r \,dr\,d\varphi = \int_0^{2\pi} \left[ \frac{r^2}{2} \right]_{r=0}^{r=1+\cos\varphi} \,d\varphi = \frac{1}{2} \int_0^{2\pi} (1+\cos\varphi)^2 \,d\varphi.\]
Expanding the integrand: $(1+\cos\varphi)^2 = 1 + 2\cos\varphi + \cos^2\varphi$. We know $\int_0^{2\pi} \cos\varphi \,d\varphi = 0$ and $\int_0^{2\pi} \cos^2\varphi \,d\varphi = \pi$. Thus,
\[\int_\Omega 1 \,dx\,dy = \frac{1}{2} (2\pi + 0 + \pi) = \frac{3\pi}{2}.\]
\end{example}

% Source: Jérôme Paschoud/Notizen/Woche 9.1 Variablenwechsel.pdf, p. 3
% Generator: Gemini 3.1 Pro (High)
\begin{example}[Combining substitution and Fubini]
\label{ex:substitution_and_fubini_rational}
Compute $\int_\Omega (x^3 y + x y^3) \,dx\,dy$ where $\Omega := \{(x,y) \in \mathbb{R}^2 \mid x,y > 0,\ 1/x \le y \le 3/x,\ 1 \le x^2 - y^2 \le 4\}$.
The integrand factors as $xy(x^2+y^2)$. The boundary inequalities are equivalent to $1 \le xy \le 3$ and $1 \le x^2-y^2 \le 4$, which strongly suggests defining $u := xy$ and $v := x^2 - y^2$.
The domain in the $(u,v)$-plane is the rectangle $R = [1,3] \times [1,4]$.
Let $\Psi(x,y) = (xy, x^2-y^2) = (u,v)$. The Jacobian of this inverse transformation is
\[\Jac\Psi(x,y) = \begin{pmatrix} y & x \\ 2x & -2y \end{pmatrix}, \qquad \det\Jac\Psi(x,y) = -2y^2 - 2x^2 = -2(x^2+y^2).\]
For the forward transformation $\Phi = \Psi^{-1}$, we have $\big|\det\Jac\Phi(u,v)\big| = \frac{1}{\lvert\det\Jac\Psi(x,y)\rvert} = \frac{1}{2(x^2+y^2)}$.
Applying the change of variables formula:
\[\int_\Omega xy(x^2+y^2) \,dx\,dy = \int_R u(x^2+y^2) \frac{1}{2(x^2+y^2)} \,du\,dv = \frac{1}{2} \int_1^4 \int_1^3 u \,du\,dv = \frac{1}{2} \cdot 3 \cdot \left[ \frac{u^2}{2} \right]_1^3 = \frac{3}{2} \cdot 4 = 6.\]
\end{example}


% Source: Corsin Nick/Class Notes/Week 8.pdf, p. 7


% --- exercises relocated here from the dissolved sheet 8 ---

% Originally: content/week-08.tex, \exercisesheet{8}, problem 8.5
% Source: exercises/Ex8_Analysis2_eng.pdf

% Source: exercises/Ex8_Analysis2_eng.pdf, pp. 1--2
\begin{exercise}[Change of variables and Jacobians \textnormal{(important)}]
\label{ex:change_of_variables_computations}
For each of the following domains and change of variables find the jacobian and the appropriate
transformed domain. There is no need to actually compute the integrals!
\begin{enumerate}[label=\textbf{(\alph*)}]
  \item $A := \{(x_1,x_2) \in \mathbb{R}^2 : x_1 > 0,\ x_2 < x_1,\ 1 < x_1^2+x_2^2 < 4\}$ and
    $x_1 = r\cos\theta$, $x_2 = r\sin\theta$. Using the change of variables formula, complete
    the dots in the following formula
    \[\int_A x_1^2\sin(x_2)\,dx_1dx_2 = \int_{\dots}\dots\,dr\,d\theta\]
  \item $B := \{(x,y) \in \mathbb{R}^2 \mid 1 < xy < 2,\ x^2 < y < 2x^2\}$ and $u := xy$,
    $v := x^2$. Using the change of variables formula, complete the dots in the following
    formula
    \[\int_B y^2e^{-xy}\,dx\,dy = \int_{\dots}\dots\,du\,dv\]
  \item $C := \{(x,y,z) \in \mathbb{R}^3 \mid 1 < z-2y < 2,\ 0 < z < 1,\ -2 < 3x+y+z < 0\}$ and
    $u := z$, $v := z-2y$, $w := 3x+y+z$. Using the change of variables formula, complete the
    dots in the following formula
    \[\int_C xyz\,dx\,dy = \int_{\dots}\dots\,du\,dv\,dw.\]
\end{enumerate}

% Supplement: Sascha Brack/Ex Sheet Hints/Ex8_Analysis2_hints.pdf, p. 1
\exhint[Sascha Brack's hint on 8.5.2]{First pin down the signs. The constraints force
$x \neq 0$ and $y \neq 0$; then $y > 0$ forces $x > 0$, so $u > 0$ and $v > 0$ throughout $B$.
That is what lets you divide by $x$ and $y$ and rearrange the inequalities without their
directions flipping.}
\exinfo{This exercise is Problem 8.5 of Problem Sheet 8, Analysis II, Spring Semester 2026.}
\exsol{sol:change_of_variables_computations}
\end{exercise}

\begin{ainote}
The integral in \textbf{8.5.3} is written $\int_C xyz\,dx\,dy$ on the official sheet, but $C$ is
a subset of $\mathbb{R}^3$ and the right-hand side carries three differentials, so this must
read $dx\,dy\,dz$. Quoted verbatim above, as the sheet is a primary source.
\end{ainote}

% An AI-Note stood here until 2026-08-09 announcing, above the exercise, that part (c)'s domain is
% the box (0,1)x(1,2)x(-2,0) and that part (b)'s second constraint becomes v^{3/2} < u < 2v^{3/2}.
% That is the answer to two of the three parts. Moved into 99-solutions.tex per the spoiler rule
% in style.md; the surviving half duplicated the AI-Note already sitting beside the solution.

\begin{aiexercise}[A substitution read off the integrand]
% Generator: Claude Opus 5 (High)
\label{ex:ai_triangle_substitution}
Let $T$ be the closed triangle with vertices $(0,0)$, $(1,0)$ and $(0,1)$, and consider
\[I := \int_T \exp\!\left(\frac{y-x}{y+x}\right)dx\,dy.\]
\begin{enumerate}[label=\textbf{(\alph*)}]
  \item Substitute $u := y-x$ and $v := y+x$. Show that the interior of $T$ corresponds to
    $\{(u,v) : 0 < v < 1,\ -v < u < v\}$, and find the factor relating $dx\,dy$ to $du\,dv$.
  \item Evaluate $I$. The inner integral is the one to do first, and it is elementary.
  \item The integrand is undefined at the origin, which lies in $T$. Why does this not affect
    the answer?
\end{enumerate}
\exsol{sol:ai_triangle_substitution}
\end{aiexercise}

% Source: old_exams/examHS24.pdf (13 February 2025), Kurzproblem 1, part 1, p. 8
% Extractor: Gemini 3.6 Flash (High)
% Correction: Claude Opus 5 (High) --- \operatorname{vol} -> \vol and J\Phi -> \Jac\Phi per the
% declared-macro rule.
\begin{exercise}[Volume of an ellipsoid via linear change of variables]
\label{ex:volume_ellipsoid_change_variables}
Let $a, b, c > 0$ and consider the solid ellipsoid
\[U := \left\{(x,y,z) \in \mathbb{R}^3 \ \middle|\ \left(\frac{x}{a}\right)^2 + \left(\frac{y}{b}\right)^2 + \left(\frac{z}{c}\right)^2 \leq 1\right\}.\]
\begin{enumerate}[label=\textbf{(\alph*)}]
  \item Define a linear change of variables $\Phi : \mathbb{R}^3 \to \mathbb{R}^3$ mapping the closed unit ball $\overline{B_1(0)} \subseteq \mathbb{R}^3$ bijectively onto $U$.
  \item Compute the Jacobian determinant $\det \Jac\Phi(u,v,w)$ of this transformation.
  \item Compute the three-dimensional Jordan volume $\vol_3(U)$ using the change of variables formula and the known volume of the unit ball.
\end{enumerate}
\exinfo{This exercise is taken from the exam of 13 February 2025 (Prof.\ Joaquim Serra), Short Problem 1, part 1, which
also asks for a labelled 3D sketch of $U$.}
\exsol{sol:volume_ellipsoid_change_variables}
\end{exercise}

\begin{ainote}
The original exam writes the ellipsoid with an equality, $(x/a)^2+(y/b)^2+(z/c)^2 = 1$, and then
asks for its volume --- but that set is the ellipsoid \emph{surface}, whose Jordan volume is $0$.
The intended object is the solid region it bounds, so the inequality is used here.
\end{ainote}

% Source: old_exams/fs2023/Prüfung.pdf (9 August 2021), Teil A, Frage 5, p. 6
% Extractor: Claude Opus 5 (High)
\begin{exercise}[A cubic shear, and the area it sweeps out]
\label{ex:cubic_shear_area}
Let $\varphi : \mathbb{R}^2 \to \mathbb{R}^2$ be given by $\varphi(x) := (x_1^3 - x_2,\; x_2 + x_1)$,
where $x = (x_1,x_2)$.
\begin{enumerate}[label=\textbf{(\alph*)}]
  \item Compute the derivative $D_x\varphi$ for every $x \in \mathbb{R}^2$.
  \item Show that $\varphi : \mathbb{R}^2 \to \varphi(\mathbb{R}^2)$ is a diffeomorphism.
  \item What is the area of the image $\varphi(Q)$ of the square $Q := [0,1] \times [0,1]$?
\end{enumerate}
\exinfo{This exercise is taken from the Analysis I \& II exam of 9 August 2021 (Prof.\ Giovanni Felder), Part A, Problem 5.}
\exsol{sol:cubic_shear_area}
\end{exercise}

% Source: old_exams/examHS24.pdf (13 February 2025), Wahr/Falsch-Fragen, Aufgabe 9, p. 5
% Extractor: Claude Opus 5 (High)
\begin{exercise}[A Gaussian second moment in polar coordinates]
\label{ex:gaussian_second_moment_polar}
For $\alpha > 0$, compute the integral
\[I_\alpha := \int_{\mathbb{R}^2} x^2 e^{-\alpha(x^2+y^2)} \, dx\,dy .\]
\exhint[Official hint]{Notice that, by symmetry, $I_\alpha = \int_{\mathbb{R}^2} y^2
e^{-\alpha(x^2+y^2)}\,dx\,dy$ as well. Compute $2I_\alpha$ using polar coordinates.}
\exinfo{This exercise is taken from the exam of 13 February 2025 (Prof.\ Joaquim Serra), where it appears as the ninth
block of true/false questions, there phrased as a series of checkpoints along this computation.}
\exsol{sol:gaussian_second_moment_polar}
\end{exercise}

% Source: old_exams/examFS24.pdf (21 August 2024), Wahr/Falsch-Fragen, Aufgabe 9, p. 5
% Extractor: Claude Opus 5 (High)
% Filed next to its February 2025 sibling above deliberately: both are one integral over R^2 in
% polar coordinates, and the pair is worth doing together because the weight differs (Gaussian
% against exponential-of-the-radius) while the method does not.
\begin{exercise}[An exponential of the radius, and how it decays in $\alpha$]
\label{ex:radial_exponential_integral}
For $\alpha > 0$, compute
\[I_\alpha := \int_{\mathbb{R}^2} e^{-\alpha\sqrt{x^2+y^2}}\,dx\,dy,\]
and then decide whether each of the following is true or false: $I_1 = \pi$; $I_2 = \pi/2$;
$\lim_{\alpha\to+\infty}\alpha I_\alpha = 0$.
\exhint[Official hint]{You may use polar coordinates and integration by parts.}
\exinfo{This exercise is taken from the exam of 21 August 2024 (Prof.\ Joaquim Serra), where it appears as the ninth block
of true/false questions. The instruction to compute $I_\alpha$ is the exam's own preamble to those
three statements.}
\exsol{sol:radial_exponential_integral}
\end{exercise}

% Source: old_exams/examHS24.pdf (13 February 2025), Wahr/Falsch-Fragen, Aufgabe 6, p. 4
% Extractor: Claude Opus 5 (High)
% Filed in this chapter rather than in 15-inverse-function-theorem, where a Gemini 3.1 Pro pass had
% put a different (and invented) block under this citation: only part (a) below is the inverse
% function theorem, and parts (b) and (c) are change of variables.
\begin{exercise}[True or False \textnormal{(the image of a small set under a $C^1$ map)}]
\label{ex:examHS24_TF_image_of_small_sets}
Consider $\Phi \in C^1(\mathbb{R}^2, \mathbb{R}^2)$ with
\[\Phi(0,0) = \begin{pmatrix} 0 \\ 0 \end{pmatrix} \quad\text{and}\quad
  \Jac\Phi(0,0) = \begin{pmatrix} 2 & 1 \\ -1 & 1 \end{pmatrix},\]
and let $X := \Phi^{-1}\big((0,0)\transp\big)$. Decide whether each of the following statements is
true or false.
\begin{enumerate}[label=\textbf{(\alph*)}]
  \item For $r > 0$ small enough, the set $\Phi(B_r(0,0))$ is open.
  \item $\vol_2\big(\Phi(B_r(0,0))\big) = 2\pi r^2 + o(r^2)$ as $r \to 0^+$.
  \item $\vol_2\big(\Phi([0,r] \times [0, ar^{1/2}])\big) = 3ar^{3/2} + o(r^{3/2})$ as
    $r \to 0^+$, where $a > 0$ is a constant.
\end{enumerate}
\exinfo{This exercise is taken from the exam of 13 February 2025 (Prof.\ Joaquim Serra), where it appears as the sixth
block of true/false questions.}
\exsol{sol:image_of_small_sets_tf}
\end{exercise}

% Source: old_exams/examHS24.pdf (13 February 2025), Wahr/Falsch-Fragen, Aufgabe 7, p. 4
% Extractor: Claude Opus 5 (High)
% Replaces a block that a Gemini 3.1 Pro pass filed under this same citation. That block asked for
% the volume of an ellipsoid spanned by three vectors of determinant 2, which is not on p. 4 and is
% nowhere on this exam; it survives, honestly relabelled, as ex:ai_spanned_ellipsoid_volume below.
\begin{exercise}[True or False \textnormal{(volumes under an affine map)}]
\label{ex:examHS24_TF_affine_volume_scaling}
Consider the map
\[\Psi : \mathbb{R}^3 \to \mathbb{R}^3, \qquad
  \Psi(x,y,z) := \big(-x + e^\pi,\; y + \alpha z,\; \alpha y + z\big),\]
where $\alpha \in \mathbb{R}$ is a parameter. Decide whether each of the following statements is
true or false.
\begin{enumerate}[label=\textbf{(\alph*)}]
  \item $\vol_3(\Psi(E)) = 0$ for all Jordan measurable $E \subseteq \mathbb{R}^3$ if and only if
    $\alpha = 1$.
  \item For $\alpha = -2$, $\vol_3(\Psi^k(E)) = 3^k \vol_3(E)$ for all Jordan measurable
    $E \subseteq \mathbb{R}^3$ and all $k \ge 0$, where $\Psi^0 := \id$ and
    $\Psi^{k+1} := \Psi \circ \Psi^k$.
  \item For $\alpha = 0$, $\vol_3(\Psi^{2025}(E)) = \vol_3(E)$ for all Jordan measurable
    $E \subseteq \mathbb{R}^3$.
  \item For all $\alpha \in (0,1)$, $\vol_3\big(\Psi^{-1}([0,1]^3)\big) = -1/(1-\alpha^2)$.
\end{enumerate}
\exinfo{This exercise is taken from the exam of 13 February 2025 (Prof.\ Joaquim Serra), where it appears as the seventh
block of true/false questions.}
\exsol{sol:affine_volume_scaling_tf}
\end{exercise}

% Source: old_exams/examFS24.pdf (21 August 2024), Wahr/Falsch-Fragen, Aufgabe 7, pp. 4--5
% Extractor: Claude Opus 5 (High)
% Worked out rather than set as an exercise, because the exercise directly above is the same
% question with a parameter in it and is strictly harder. What this block adds is the composition
% and inverse rules read off one determinant, which the February 2025 version states only for
% iterates of a single map.
\begin{example}[One determinant, three volume statements]
\label{ex:affine_volume_one_determinant}
% Generator: Claude Opus 5 (High) --- the three statements are the examiner's; the working is
% written for these notes.
Consider $\Psi : \mathbb{R}^3 \to \mathbb{R}^3$, $\Psi(x,y,z) := (z-1,\ x+2y,\ x-y)$. Its linear
part has matrix
\[A = \begin{pmatrix} 0 & 0 & 1 \\ 1 & 2 & 0 \\ 1 & -1 & 0\end{pmatrix},
  \qquad \det A = 1\cdot\det\begin{pmatrix} 1 & 2 \\ 1 & -1\end{pmatrix} = -3,\]
expanding along the first row. The translation by $(-1,0,0)$ does not affect volumes, so
$\vol_3(\Psi(E)) = 3\vol_3(E)$ for every Jordan measurable $E$, by
\cref{thm:change_of_variables}. Everything else follows from the factor $3$.
\begin{enumerate}[label=\textbf{(\alph*)}]
  \item \textbf{$\vol_3(\Psi^{-1}(E)) = \vol_3(E)$ for all Jordan measurable $E$: false.} The
    inverse scales by $\lvert\det A^{-1}\rvert = 1/3$, so $\vol_3(\Psi^{-1}(E)) = \vol_3(E)/3$.
  \item \textbf{$\vol_3(\Psi(E)) = 2\vol_3(E)$: false.} The factor is $3$, not $2$.
  \item \textbf{$\vol_3(\Psi(\Psi(E))) = 9\vol_3(E)$: true.} Applying \textbf{(b)}'s correct factor
    twice gives $3 \cdot 3 = 9$, which is $\lvert\det A^2\rvert = \lvert\det A\rvert^2$.
\end{enumerate}
Only \textbf{(b)} requires the determinant to be computed. Items \textbf{(a)} and \textbf{(c)} then
follow from multiplicativity of the determinant alone, without recomputing anything, which is the
point of the block.
\exhint[Official hint]{To answer these three, it can be useful to compute the Jacobian determinant
of $\Psi$ first.}
\exinfo{This example is taken from the exam of 21 August 2024 (Prof.\ Joaquim Serra), where it appears as the seventh
block of true/false questions.}
\end{example}

% Source: old_exams/examHS24.pdf (13 February 2025), Wahr/Falsch-Fragen, Aufgabe 8, pp. 4--5
% Extractor: Gemini 3.1 Pro (High)
\begin{exercise}[True or False \textnormal{(spherical coordinates)}]
\label{ex:examHS24_TF_spherical}
Consider the polar (spherical) coordinates $\Psi \colon (0, 2\pi) \times (0, \pi) \times (0, \infty)
\to \mathbb{R}^3$ defined by
\[ \Psi(\phi, \theta, r) = (r \cos \phi \sin \theta, r \sin \phi \sin \theta, r \cos \theta). \]
Decide whether each of the following statements is true or false.
\begin{enumerate}[label=\textbf{(\alph*)}]
  \item $\Psi$ is injective.
  \item $\Psi$ is surjective.
  \item The Jacobian determinant $\det \Jac\Psi(\theta, r)$ is given by $r^2 \sin \theta$.
  \item $\vol_3(B_R \setminus B_r) = \pi(R^3 - r^3)$.
  \item $\vol_3(B_R) = 2\pi \int_0^\pi \int_0^R r^2 \sin \theta \, dr \, d\theta$.
  \item $\lim_{R \to 1^+} \frac{\vol_3(B_R \setminus B_1)}{R - 1} = 4\pi$.
\end{enumerate}
\exinfo{This exercise is taken from the exam of 13 February 2025 (Prof.\ Joaquim Serra), where it appears as the eighth
block of true/false questions.}
\exsol{sol:spherical_coordinates_integral_tf}
\end{exercise}

% Source: old_exams/examHS24.pdf (13 February 2025), Kästchenaufgabe 2, p. 7
% Extractor: Gemini 3.1 Pro (High)
\begin{exercise}[Box question: Linear map scaling volume]
\label{ex:examHS24_box_linear_vol}
Give an explicit example of a linear map $\Phi \colon \mathbb{R}^2 \to \mathbb{R}^2$ such that
\[ \vol_2(\Phi(E)) = 2 \vol_2(E) \quad \text{for all } E \subset \mathbb{R}^2 \text{ Jordan measurable.} \]
\exinfo{This exercise is taken from the exam of 13 February 2025 (Prof.\ Joaquim Serra), Box Problem 2, where only the
final answer was graded.}
\exsol{sol:linear_map_volume_scaling_box}
\end{exercise}

% Source: old_exams/examFS24.pdf (21 August 2024), Kästchenaufgabe 2, p. 7
% Extractor: Claude Opus 5 (High)
% The August 2024 counterpart of the box problem directly above: same one-line format, factor 1
% instead of factor 2, and the added clause "different from the identity", which is what stops the
% obvious answer and makes the problem worth having next to its sibling.
\begin{exercise}[Box question: a linear map that preserves every volume]
\label{ex:box_linear_volume_preserving}
Give an explicit example of a linear map $\Phi : \mathbb{R}^2 \to \mathbb{R}^2$, \textbf{different
from the identity}, such that
\[\vol_2(\Phi(E)) = \vol_2(E) \quad \text{for all Jordan measurable } E \subseteq \mathbb{R}^2 .\]
\exinfo{This exercise is taken from the exam of 21 August 2024 (Prof.\ Joaquim Serra), Box Problem 2, where only the final
answer was graded.}
\exsol{sol:box_linear_volume_preserving}
\end{exercise}

% Correction: Claude Opus 5 (High) --- provenance repair. A Gemini 3.1 Pro pass filed the statement
% below as an exam problem, citing old_exams/examHS24.pdf, Aufgabe 7, p. 4. That page carries
% ex:examHS24_TF_affine_volume_scaling instead, and nothing resembling this problem is anywhere on
% that exam. The mathematics is sound, so the block is kept as what it actually is: an exercise
% authored for these notes. Its % Source:, % Extractor: and \exinfo lines have been removed.
\begin{aiexercise}[The volume of a spanned ellipsoid]
\label{ex:ai_spanned_ellipsoid_volume}
% Generator: Gemini 3.1 Pro (High)
Let $v_1, v_2, v_3 \in \mathbb{R}^3$ be vectors with $\det(v_1, v_2, v_3) = 2$, and consider the
region
\[ A := \big\{ t_1 v_1 + t_2 v_2 + t_3 v_3 \ \big|\ (t_1, t_2, t_3) \in \mathbb{R}^3
  \text{ with } t_1^2 + t_2^2 + t_3^2 \le 1 \big\}. \]
Decide whether each of the following statements is true or false.
\begin{enumerate}[label=\textbf{(\alph*)}]
  \item $\vol_3(A) = \frac{8\pi}{3}$.
  \item $A$ is Jordan measurable.
  \item The volume of $A$ depends on the choice of the vectors $v_1, v_2, v_3$.
\end{enumerate}
\exsol{sol:ai_spanned_ellipsoid_volume}
\end{aiexercise}

% Source: old_exams/HS18.pdf (Analysis I & II, winter 2018/2019, Prof. Manfred Einsiedler),
% Aufgabe 13, p. 4
% Extractor: Claude Opus 5 (High)
% This is the proof of thm:jordan_measure_linear_transformation, which chapter 19 states without
% one. Filed here rather than beside the theorem because the exam frames it as the linear case of
% the substitution rule, which is this chapter's subject.
\begin{exercise}[The substitution rule in the linear case]
\label{ex:linear_substitution_elementary_matrices}
Let $A$ be an invertible $n\times n$ matrix. \Cref{thm:jordan_measure_linear_transformation}
asserts that for every Jordan measurable $B \subseteq \mathbb{R}^n$ the image $A(B)$ is again
Jordan measurable, with
\begin{equation}
\label{eq:linear_substitution_rule}
\vol\big(A(B)\big) = \lvert\det A\rvert \, \vol(B).
\end{equation}
This exercise proves it.
\begin{enumerate}[label=\textbf{(\alph*)}]
  \item Prove \eqref{eq:linear_substitution_rule} directly in the case that $A$ is a diagonal
    matrix.
  \item Explain in a few sentences why it now suffices to prove \eqref{eq:linear_substitution_rule}
    for permutation matrices and for shear matrices.
  \item Prove \eqref{eq:linear_substitution_rule} when $A$ is a shear matrix. You may assume for
    simplicity that $B = Q$ is a closed axis-parallel box.
\end{enumerate}
\exinfo{This exercise is taken from the Analysis I \& II examination of winter 2018/2019
(Prof.\ Manfred Einsiedler), Problem 13.}
\exsol{sol:linear_substitution_elementary_matrices}
\end{exercise}

% Source: old_exams/HS17.pdf (Analysis I & II, winter 2017/2018, Prof. Manfred Einsiedler),
% Aufgabe 13, pp. 3--4
% Extractor: Claude Opus 5 (High)
% Kept as a worked example, not an exercise: this is the one genuinely hard lemma behind
% thm:change_of_variables, which this document states without proof, and it is well beyond what
% any exercise here asks. The exam splits it into four graded parts; those are preserved as the
% four steps below.
% Notation normalised to the house convention: the exam writes ||.||_2 for the Euclidean norm,
% which is single bars here (style.md, "Norms"). ||.||_infty keeps its double bars.
\begin{example}[The covering lemma behind the change of variables formula]
\label{ex:cube_covering_lemma_substitution}
\Cref{thm:change_of_variables} is stated in this chapter without proof. The step that makes the
proof hard is not the Jacobian factor but a prior, purely geometric one: a diffeomorphism that is
close to the identity cannot let a cube shrink away. Made precise, that is the following.

\textbf{Claim.} Let $\ell > 0$, let $Q_0 := [-\ell,\ell]^n \subseteq \mathbb{R}^n$, let
$X \supseteq Q_0$ be open and let $\Phi : X \to \Phi(X) \subseteq \mathbb{R}^n$ be a
diffeomorphism with $\Phi(0) = 0$ and $D_0\Phi = I_n$. Suppose $\sigma > 0$ is such that
\[\big\lvert (D_x\Phi - I_n)v \big\rvert \le \sigma \lvert v\rvert
\qquad \text{for all } x \in Q_0,\ v \in \mathbb{R}^n,\]
and assume $s := \sigma\sqrt{n} < 1$. Then
\[(1-s)\,Q_0 \subseteq \Phi(Q_0).\]

Fix $y \in (1-s)Q_0$ and define $F_y : Q_0 \to \mathbb{R}^n$ by
$F_y(x) := y - \big(\Phi(x) - x\big)$.

\begin{enumerate}[label=\textbf{(\alph*)}]
  \item \textbf{Why a fixed point is what we want.} By construction $F_y(x) = x$ holds if and only
    if $y - \Phi(x) + x = x$, that is, if and only if $\Phi(x) = y$. So a fixed point of $F_y$ in
    $Q_0$ is exactly a preimage of $y$ under $\Phi$ lying in $Q_0$, which is what
    $y \in \Phi(Q_0)$ means.

  \item \textbf{The displacement $\Phi - \id$ is a contraction.} Write $\Psi(x) := \Phi(x) - x$, so
    that $D_x\Psi = D_x\Phi - I_n$ and the hypothesis reads
    $\lvert D_x\Psi(v)\rvert \le \sigma\lvert v\rvert$ on $Q_0$, that is,
    $\lVert D_x\Psi\rVert_{\mathrm{op}} \le \sigma$ there. Since $Q_0$ is convex, the segment from
    $x_1$ to $x_2$ stays inside it, and the mean value inequality
    (\cref{thm:mean_value_inequality}) gives
    \[\big\lvert \Phi(x_2)-x_2 - \big(\Phi(x_1)-x_1\big)\big\rvert
      = \lvert \Psi(x_2)-\Psi(x_1)\rvert \le \sigma \lvert x_2-x_1\rvert\]
    for all $x_1,x_2 \in Q_0$.

  \item \textbf{$F_y$ maps $Q_0$ into itself.} Applying \textbf{(b)} with $x_1 := 0$, and using
    $\Psi(0) = \Phi(0)-0 = 0$, we get $\lvert \Psi(x)\rvert \le \sigma\lvert x\rvert$ for
    $x \in Q_0$. For $x \in Q_0$ we have $\lVert x\rVert_\infty \le \ell$ and hence
    $\lvert x\rvert \le \sqrt{n}\,\ell$, so
    \[\lVert \Psi(x)\rVert_\infty \le \lvert \Psi(x)\rvert \le \sigma\sqrt{n}\,\ell = s\ell.\]
    Since $y \in (1-s)Q_0$ means $\lVert y\rVert_\infty \le (1-s)\ell$, the triangle inequality
    gives
    \[\lVert F_y(x)\rVert_\infty \le \lVert y\rVert_\infty + \lVert \Psi(x)\rVert_\infty
      \le (1-s)\ell + s\ell = \ell,\]
    that is, $F_y(x) \in Q_0$. So $F_y : Q_0 \to Q_0$.

  \item \textbf{Conclusion.} $Q_0$ is a closed subset of $\mathbb{R}^n$, hence complete, and by
    \textbf{(b)} the map $F_y$ satisfies
    $\lvert F_y(x_2)-F_y(x_1)\rvert = \lvert \Psi(x_2)-\Psi(x_1)\rvert \le
    \sigma\lvert x_2-x_1\rvert$ with $\sigma \le s < 1$. It is therefore a contraction of the
    complete metric space $Q_0$ into itself, and the Banach fixed point theorem
    (\cref{thm:banach_fixed_point}) supplies a fixed point $x \in Q_0$. By \textbf{(a)},
    $\Phi(x) = y$, so $y \in \Phi(Q_0)$. As $y \in (1-s)Q_0$ was arbitrary,
    $(1-s)Q_0 \subseteq \Phi(Q_0)$.
\end{enumerate}

\begin{remark}[What the lemma is doing in the proof of the change of variables formula]
The formula is proved by cutting the domain into tiny cubes and comparing $\vol(\Phi(Q))$ with
$\lvert\det D\Phi\rvert\vol(Q)$ on each. The upper bound is the easy half: $\Phi$ is Lipschitz on a
compact piece, so $\Phi(Q)$ cannot be much bigger than $Q$. The lower bound is where a naive
argument stalls, because nothing about a Lipschitz map prevents an image from collapsing. This
lemma is that lower bound. First normalise so that $D_0\Phi = I_n$, which is always possible by
composing with a linear map, the case already settled in
\cref{ex:linear_substitution_elementary_matrices}. The lemma then says the image of a cube still
contains a cube shrunk only by the factor $1-s$, and $s \to 0$ as the cubes shrink, because
$D\Phi$ is continuous. Letting the two bounds close is the proof.

Notice which two theorems do the work: the mean value inequality on a convex set in \textbf{(b)},
and the Banach fixed point theorem in \textbf{(d)}. The same pair proves the inverse function
theorem (\cref{thm:inverse_function_theorem}), and that is not a coincidence. This claim is a
quantitative version of the statement that $\Phi$ is locally surjective near a point where its
differential is invertible.
\end{remark}
\exinfo{This example is taken from the Analysis I \& II examination of winter 2017/2018
(Prof.\ Manfred Einsiedler), Problem 13, where the four steps above are set as four graded parts.
The framing and the closing remark are added here.}
\end{example}

